\section{Điều khiển tắc nghẽn và Head-of-Line blocking}
\label{sec:cc-hol}

Chương này trình bày hai yêu cầu quan trọng của đề: điều khiển tắc nghẽn, mà nhóm chọn cài đặt bằng thuật
toán CUBIC, và loại bỏ Head-of-Line blocking giữa tin nhắn và file. Với mỗi phần, nhóm đều có phép đo để kiểm
chứng.

\subsection{CUBIC}
\label{subsec:cubic}
Điều khiển luồng (\cref{subsec:dieu-toc}) giúp bảo vệ \emph{bên nhận} khỏi bị gửi quá nhiều, còn điều khiển tắc
nghẽn giúp bảo vệ \emph{mạng}. Bên gửi giữ một cửa sổ tắc nghẽn gọi là \textit{cwnd} (\emph{congestion window}),
là số byte tối đa được phép đang trên đường truyền mà chưa được ACK. Giá trị cwnd được điều chỉnh dựa trên tín
hiệu mất gói. Nhóm cài đặt CUBIC theo RFC~9438 \cite{rfc9438}; đây cũng là thuật toán mặc định trong Linux. Các
quy tắc chính như sau:

\begin{itemize}
  \item \textbf{Slow start:} lúc mới kết nối, cwnd bằng 10 gói (12 KB). Mỗi byte được ACK thì cwnd tăng thêm
        1 byte, nghĩa là cwnd tăng gấp đôi sau mỗi RTT.
  \item \textbf{Khi mất gói:} ghi nhớ $W_\text{max} \leftarrow \mathit{cwnd}$, giảm
        $\mathit{cwnd} \leftarrow \beta \cdot \mathit{cwnd}$ với $\beta = 0{,}7$, và thoát khỏi slow start.
        Mỗi đợt nghẽn chỉ được giảm một lần: nếu các gói gửi trước lần giảm gần nhất bị mất thêm thì không giảm
        tiếp nữa. Nếu không có quy tắc này, cwnd sẽ bị chia liên tục chỉ vì cùng một lần nghẽn.
  \item \textbf{Tăng lại} theo hàm bậc 3 của thời gian $t$ tính từ lần mất gần nhất (đơn vị là gói):
        \begin{equation}
          W_\text{cubic}(t) = C\,(t-K)^3 + W_\text{max}, \qquad
          K = \sqrt[3]{\frac{W_\text{max}(1-\beta)}{C}}, \qquad C = 0{,}4
          \label{eq:cubic}
        \end{equation}
        Đường cong này tăng nhanh khi cwnd còn xa $W_\text{max}$, sau đó chậm dần và đi ngang quanh
        $W_\text{max}$ (là mức đã từng gây nghẽn), rồi lại tăng nhanh để thử xem mạng có còn chỗ hay không. Vì
        cwnd tăng theo thời gian thực chứ không theo số ACK nhận được, nên kết nối có RTT dài cũng hồi phục nhanh
        như kết nối có RTT ngắn.
  \item \textbf{Vùng Reno-friendly:} song song với $W_\text{cubic}$, nhóm tính thêm $W_\text{est}$ là cửa sổ mà
        thuật toán Reno (thuật toán cũ của TCP) sẽ có nếu chạy trong cùng điều kiện. Reno tăng
        $\alpha = 3\frac{1-\beta}{1+\beta} \approx 0{,}53$ gói mỗi RTT. Sau đó lấy
        $\mathit{cwnd} = \max(W_\text{cubic}, W_\text{est})$. Lý do là khi RTT ngắn hoặc cửa sổ nhỏ, Reno lại tăng
        nhanh hơn đường bậc 3, mà CUBIC thì không được phép chạy tệ hơn Reno.
  \item \textbf{Giới hạn:} cwnd có sàn là 2 gói, vì nếu cwnd dưới 1 gói thì không gửi được gì. Trần của cwnd là
        1 MB, bằng kích thước cửa sổ nhận, vì có tăng hơn nữa thì cũng không gửi được nhiều hơn.
\end{itemize}

\begin{lstlisting}[caption={Tăng cwnd khi nhận ACK (\fname{cubic.go}, rút gọn)}, label=lst:cubic]
func (c *cubic) OnAck(acked int, now time.Time) {
	a := float64(acked)
	if c.cwnd < c.ssthresh { // slow start: doubles every RTT
		c.cwnd += a
		return
	}
	if c.epoch.IsZero() { // first ACK after a loss: start a new cubic epoch
		c.epoch, c.wEst = now, c.cwnd
	}
	c.wEst += cubicAlpha * mss * a / c.cwnd // Reno estimate: +alpha packets per RTT
	cub := c.wCubic(now.Sub(c.epoch).Seconds())
	if cub < c.wEst {
		c.cwnd = math.Max(c.cwnd, c.wEst) // Reno-friendly region
		return
	}
	target := math.Min(cub, 1.5*c.cwnd) // never more than 1.5x per RTT
	if target > c.cwnd {
		c.cwnd += (target - c.cwnd) * a / c.cwnd
	}
}
\end{lstlisting}

\paragraph{Kiểm chứng.} Để quan sát hành vi của CUBIC, nhóm dùng bộ giả lập mạng (\cref{subsec:gia-lap}) tạo
một nút cổ chai có băng thông 2 MB/s, trễ 20 ms mỗi chiều và bộ đệm 64 KB. Nhóm cố ý \textbf{không bật mất gói
ngẫu nhiên}, để mọi lần mất gói đều do bộ đệm bị tràn, tức là do nghẽn thật. Nhóm truyền 60 MB và ghi lại giá
trị cwnd ở mỗi nhịp đồng hồ (test \fname{TestLongCCGraph} với biến môi trường \texttt{BTL\_CCLOG}). Kết quả được
vẽ trong \cref{fig:cwnd,fig:cwnd-zoom}.

\begin{figure}[H]
  \centering
  \cachedtikz{cwnd}{%
  \begin{tikzpicture}
    \begin{axis}[
        width=0.97\linewidth, height=6.5cm,
        xlabel={Thời gian (s)}, ylabel={cwnd (KB)},
        xmin=0, xmax=31.5, ymin=0, ymax=310,
        grid=major, grid style={gray!20},
        legend pos=north east, legend cell align=left, legend style={font=\small},
        tick label style={font=\small}, label style={font=\small},
      ]
      \addplot[blue!70!black, thick, mark=none] table[header=false, x index=0, y index=1] {data/cwnd.dat};
      \addlegendentry{cwnd}
      \addplot[only marks, mark=x, mark size=3pt, red!80!black, thick]
        table[header=false, x index=0, y index=1] {data/loss.dat};
      \addlegendentry{mất gói (cwnd ngay trước khi giảm)}
      \addplot[gray, dashed, domain=0:31.5] {148};
      \addlegendentry{BDP + bộ đệm $\approx$ 148 KB}
    \end{axis}
  \end{tikzpicture}}
  \caption{cwnd theo thời gian: CUBIC qua nút cổ chai 2 MB/s, RTT 40 ms, bộ đệm 64 KB}
  \label{fig:cwnd}
\end{figure}

\begin{figure}[H]
  \centering
  \cachedtikz{cwnd-zoom}{%
  \begin{tikzpicture}
    \begin{axis}[
        width=0.97\linewidth, height=6cm,
        xlabel={Thời gian (s)}, ylabel={cwnd (KB)},
        xmin=3, xmax=16, ymin=95, ymax=155,
        grid=major, grid style={gray!20},
        tick label style={font=\small}, label style={font=\small},
      ]
      \addplot[blue!70!black, thick, mark=none] table[header=false, x index=0, y index=1] {data/cwnd.dat};
      \addplot[only marks, mark=x, mark size=3pt, red!80!black, thick]
        table[header=false, x index=0, y index=1] {data/loss.dat};
    \end{axis}
  \end{tikzpicture}}
  \caption{Phóng to \cref{fig:cwnd} trong đoạn 3--16 s: đường bậc 3 ``chữ S nằm ngang'' giữa hai lần mất}
  \label{fig:cwnd-zoom}
\end{figure}

Từ hai đồ thị, nhóm rút ra một số nhận xét:
\begin{itemize}
  \item Ở giai đoạn slow start, cwnd tăng vọt lên khoảng 297 KB, sau đó gặp ba lần mất liên tiếp trong 0,6 s đầu
        (297 $\to$ 210 $\to$ 148 $\to$ 105 KB). Lý do là cwnd đã vượt quá xa sức chứa của đường truyền nên phải
        giảm vài lần mới xuống dưới mức gây nghẽn.
  \item Sau giai đoạn đầu, mất gói luôn xảy ra khi cwnd vào khoảng 147 KB. Con số này khớp với sức chứa của
        đường truyền, bằng BDP (\emph{bandwidth-delay product}, tích băng thông và độ trễ)
        $= 2\,\text{MB/s} \times 42\,\text{ms} \approx 84$ KB cộng thêm bộ đệm 64 KB. Sau mỗi lần mất, cwnd còn
        khoảng 103 KB, đúng bằng $0{,}7 \times 147$ KB.
  \item Giữa hai lần mất, cwnd tăng nhanh rồi đi ngang khi tới gần $W_\text{max}$ (\cref{fig:cwnd-zoom}), đúng
        với hình dạng của công thức~\eqref{eq:cubic}. Chu kỳ đo được khoảng 3,3--4,8 s, trung bình 4 s. Theo lý
        thuyết, với $W_\text{max} \approx 126$ gói thì $K = \sqrt[3]{126 \times 0{,}3 / 0{,}4} \approx 4{,}5$ s,
        khá gần với kết quả đo.
  \item CUBIC đạt \textbf{2,01 MB/s}, tức là dùng hết băng thông của nút cổ chai. Để so sánh, nhóm chạy thử với
        cửa sổ cố định 64 KB (\texttt{-cc fixed}) trên cùng đường truyền thì chỉ đạt 1,47 MB/s. Nguyên nhân là
        64 KB nhỏ hơn BDP nên đường truyền không bao giờ được dùng hết.
\end{itemize}

\subsection{Giới hạn của CUBIC khi mất gói ngẫu nhiên}
\label{subsec:cubic-gioi-han}
Trong quá trình đo, nhóm nhận thấy CUBIC hoạt động kém khi mạng mất gói ngẫu nhiên. Giống như mọi thuật toán
dựa trên mất gói, CUBIC coi \emph{mọi} lần mất gói là dấu hiệu nghẽn. Vì vậy khi mạng mất gói ngẫu nhiên (ví dụ
sóng wifi yếu, hoặc giả lập bằng \texttt{-loss 0.1}) mà thực ra không hề nghẽn, cwnd vẫn liên tục bị giảm và gần
như luôn nằm sát mức sàn. Công thức Mathis \cite{mathis} cho thấy throughput của các thuật toán loại này tỉ lệ
nghịch với căn bậc hai của tỉ lệ mất gói $p$:
\begin{equation}
  \text{throughput} \approx \frac{\mathit{MSS}}{\mathit{RTT}} \cdot \frac{1{,}22}{\sqrt{p}}
\end{equation}
Kết quả đo thực tế cũng đúng như vậy. Với tỉ lệ mất 10\% và jitter 20 ms, CUBIC mất 48,3 s để truyền 5 MB
(0,11 MB/s), trong khi cửa sổ cố định 64 KB truyền 50 MB đạt 1,69 MB/s (\cref{tab:do-transport}). Do đó nhóm
giữ lại tùy chọn \texttt{-cc fixed} để demo độ tin cậy của giao thức khi mất gói nặng. Cách giải quyết triệt để
hơn là dùng BBR \cite{bbr}, thuật toán ước lượng băng thông và RTT tối thiểu thay vì dựa vào mất gói. Nhóm để
đây là hướng phát triển (\cref{sec:ket-luan}).

\subsection{Không có Head-of-Line blocking}
\label{sec:hol}
\emph{Head-of-Line blocking} (HoL) là hiện tượng dữ liệu đứng sau phải chờ dữ liệu đứng trước, dù hai loại dữ
liệu này không liên quan gì tới nhau. Ví dụ với TCP, nếu một gói của file bị mất thì \emph{mọi} byte phía sau
nó, kể cả tin chat, đều phải chờ gói đó được gửi lại. Đề bài yêu cầu loại bỏ hiện tượng này, và nhóm giải quyết ở
cả hai phía:

\begin{enumerate}
  \item \textbf{Phía nhận: mỗi stream được ghép độc lập.} Mỗi tin nhắn và mỗi file là một stream riêng, có
        \fname{recvStream} riêng (\cref{subsec:ghep}). Khi một gói của file bị mất thì chỉ stream của file phải
        chờ, còn stream chat nhận đủ dữ liệu là được giao ngay.
  \item \textbf{Phía gửi: hàng đợi ưu tiên.} Tuy nhiên, chỉ tách stream thì chưa đủ:
        nếu tin chat phải xếp hàng sau hàng trăm frame file đang chờ gửi thì nó vẫn bị chậm. Vì vậy hàm
        \fname{takeFrames} (\cref{lst:takeframes}) lấy frame theo thứ tự \texttt{retransQ} $\to$
        \texttt{urgentQ} $\to$ \texttt{bulkQ}, nhờ đó tin chat được chen lên trước toàn bộ dữ liệu file đang
        chờ. Dữ liệu cần gửi lại được ưu tiên cao nhất để các stream đang thiếu dữ liệu được lấp đầy sớm.
\end{enumerate}

\begin{lstlisting}[caption={Lấy frame cho một gói theo thứ tự ưu tiên (\fname{conn.go}, rút gọn)},
  label=lst:takeframes]
func (c *Conn) takeFrames() []*protocol.StreamFrame {
	room := protocol.MaxFramePayload - ackFrameMax // leave room for an ACK
	var out []*protocol.StreamFrame
	for _, q := range []*[]*protocol.StreamFrame{&c.retransQ, &c.urgentQ, &c.bulkQ} {
		for len(*q) > 0 {
			f := (*q)[0]
			size := streamFrameHdr + len(f.Data)
			if size > room {
				return out // packet is full
			}
			room -= size
			out = append(out, f)
			*q = (*q)[1:]
		}
	}
	return out
}
\end{lstlisting}

\paragraph{Kiểm chứng} (\fname{TestLongHoL}). Nhóm gửi một file 6 MB qua nút cổ chai 1 MB/s, trễ 20 ms mỗi
chiều, bộ đệm 64 KB và tỉ lệ mất gói 1\%. Trong lúc gửi file, cứ mỗi 100 ms nhóm gửi một tin chat và đo thời gian
từ lúc gửi tới lúc bên kia nhận được. Nhóm so sánh thiết kế hiện tại với một phiên bản cố tình cho tin chat xếp
chung hàng với file, kết quả ở \cref{tab:hol}.

\begin{table}[H]
  \centering
  \caption{Độ trễ tin chat khi đang truyền file (đo ngày 12/11/2026)}
  \label{tab:hol}
  \small
  \begin{tabularx}{\linewidth}{Yrrrr}
    \toprule
    \textbf{Tin chat đi bằng} & \textbf{Trung bình} & \textbf{p50} & \textbf{p95} & \textbf{max} \\
    \midrule
    Không có file (mốc: chỉ có trễ đường truyền) & 21,4 ms & 21,4 ms & 21,8 ms & 22,0 ms \\
    \textbf{Hàng urgent (thiết kế hiện tại)} & \textbf{27,3 ms} & \textbf{23,6 ms} & \textbf{41,3 ms} &
    \textbf{56,3 ms} \\
    Chung hàng với file (giả lập không ưu tiên) & 989,7 ms & 972,5 ms & 1470 ms & 1723 ms \\
    \bottomrule
  \end{tabularx}
\end{table}

Kết quả cho thấy khi đang gửi file, tin chat chỉ chậm hơn lúc rảnh vài ms (p50 là 23,6 ms so với 21,4 ms).
Ngược lại, nếu không có hàng đợi ưu tiên thì tin chat bị chậm gần 1 giây, tức là chậm hơn khoảng 40 lần, vì nó
phải chờ hết phần dữ liệu file đang xếp hàng phía trước và đang nằm trong bộ đệm của nút cổ chai.
